% Section: 谱维流与时空拓扑
\section{Spectral Dimension Flow and Spacetime Topology}
\label{sec:spectral_topology}

\subsection{Fundamental Distinction: Topological vs Spectral Dimension}

A crucial clarification is necessary regarding the nature of dimension in our framework. We distinguish between:

\begin{definition}[Topological Dimension]
The topological dimension $d_t$ is the number of independent coordinates needed to specify a point in the observable manifold. In our theory:
\begin{itemize}
    \item Observable spacetime: $d_t = 4$ (fixed, immutable)
\end{itemize}
The topological dimension never changes for external observers; it is a property of the manifold structure itself. The internal space is not a separate manifold with its own topological dimension, but rather the same physical system described by $Cl(1,9)$ algebra at high energies, where 1 time and 9 spatial degrees of freedom become accessible.
\end{definition}

\begin{definition}[Spectral Dimension]
The spectral dimension $d_s$ is the effective dimension "experienced" by energy propagating through spacetime. It depends on the energy scale $E$ and the local torsion field strength:
\begin{equation}
    d_s(E, \tau) = 4 + \frac{6\tau(E)}{1 + \tau(E)}
    \label{eq:spectral_dimension}
\end{equation}
where $\tau(E)$ is the energy-dependent torsion field strength.
\end{definition}

\begin{remark}[Physical Interpretation]
The spectral dimension describes how many degrees of freedom are accessible to energy at a given scale. At low energies ($E \ll E_{EW}$), $d_s \approx 4$ and energy is constrained to the 4D external spacetime. At high energies ($E \gg E_{GUT}$), $d_s \to 10$ and energy can access the full 10D structure.
\end{remark}

\subsection{The Kaleidoscopic Projection Mechanism}

The relationship between external 4D spacetime and internal 10D space is established through the kaleidoscopic projection:

\begin{equation}
    P_{kaleido}: \mathbb{R}^{10} \to \mathbb{R}^4
    \label{eq:kaleidoscopic_projection}
\end{equation}

This projection is characterized by the modulation factor:
\begin{equation}
    K(\theta) = |\sin(n_{sym} \cdot \theta)|
    \label{eq:kaleidoscope_factor}
\end{equation}
where $n_{sym} = 4$ corresponds to quaternion symmetry.

The projection angles for different force carriers are:
\begin{equation}
    \theta_n = \frac{\pi}{4} \cdot \frac{n}{n_{max}} = \frac{\pi}{4} \cdot \frac{n}{3}
    \label{eq:projection_angles}
\end{equation}

\begin{table}[htbp]
\centering
\caption{Geometric Origins of Fundamental Forces}
\begin{tabular}{@{}ccccc@{}}
\toprule
Force & Order $n$ & Internal Dim. & Gauge Group & Projection Angle \\
\midrule
Electromagnetic & 1 & 2D & $U(1)_Y$ & $15°$ (shallow) \\
Weak & 2 & 4D & $SU(2)_L$ & $30°$ (intermediate) \\
Strong & 3 & 6D & $SU(3)_C$ & $45°$ (deep) \\
Gravity & 0 & 6D & Diffeomorphism & Background \\
\bottomrule
\end{tabular}
\label{tab:force_origins}
\end{table}

\subsection{Energy State Transitions Across the Horizon}

When matter crosses the event horizon of a black hole, it undergoes a fundamental state transition:

\begin{axiom}[Energy State Duality]
Energy exists in two fundamental states:
\begin{enumerate}
    \item \textbf{4D Constrained State}: Energy is confined to 4D external spacetime with spectral dimension $d_s \approx 4$ (1 time + 3 space dimensions), topological dimension $d_t = 4$.
    \item \textbf{10D Expanded State}: Energy is released into the 10D internal space with spectral dimension $d_s \to 10$ (1 time + 9 space dimensions), while the topological dimension for external observers remains $d_t = 4$.
\end{enumerate}
\end{axiom}

The transition occurs at the event horizon $r = r_s$:

\begin{equation}
    \rho_{4D}(r) = \rho_{total} \cdot f_{4D}(r), \quad \rho_{10D}(r_{int}) = \rho_{total} \cdot f_{10D}(r_{int})
\end{equation}

with the constraint:
\begin{equation}
    f_{4D}(r) + f_{10D}(r_{int}) = 1
\end{equation}

At the horizon:
\begin{equation}
    f_{4D}(r_s) = f_{10D}(r_s) = 0.5
    \label{eq:horizon_energy_equipartition}
\end{equation}

This represents critical coupling where energy is equally partitioned between external and internal spaces.

\subsection{Inward Fractal Expansion}

\begin{theorem}[Inward Fractal Expansion]
As matter falls toward the black hole center ($r \to 0$ in external coordinates), the internal radial coordinate expands as:
\begin{equation}
    r_{int} = r_s \cdot \left(\frac{r_s}{r}\right)^{\tau_0}
    \label{eq:internal_expansion}
\end{equation}
where $\tau_0 \approx 0.005$ is the characteristic torsion strength.
\end{theorem}

\begin{corollary}[Infinite Internal Space]
As $r \to 0$, $r_{int} \to \infty$. The "singularity" in external coordinates corresponds to infinite expansion in internal space.
\end{corollary}

\begin{proof}
From equation \eqref{eq:internal_expansion}, taking the limit:
\begin{equation}
    \lim_{r \to 0} r_{int} = r_s \cdot \lim_{r \to 0} \left(\frac{r_s}{r}\right)^{\tau_0} = \infty
\end{equation}
for any positive $\tau_0$.
\end{proof}

The fractal dimension varies with depth:
\begin{equation}
    D_f(r) = 4 + \frac{6\tau(r)}{1 + \tau(r)}, \quad \tau(r) = \tau_0 \cdot \frac{r_s}{r}
\end{equation}

\begin{table}[htbp]
\centering
\caption{Dimensional Evolution Across Black Hole Horizon}
\begin{tabular}{@{}cccccc@{}}
\toprule
Region & $r/r_s$ & $d_s$ & $d_t$ & $n_{time}$ & $n_{space}$ \\
\midrule
Far exterior & $\gg 1$ & $\sim 4$ & 4 & 1 & 3 \\
Near horizon & $\sim 1$ & $4 \to 7$ & 4 & 1 & $3 \to 6$ \\
At horizon & $= 1$ & $\sim 7$ & 4 & 1 & 6 \\
Interior & $\ll 1$ & $7 \to 10$ & 4 & 1 & $6 \to 9$ \\
\bottomrule
\end{tabular}
\label{tab:dimensional_evolution}
\end{table}

\begin{remark}[Fixed Time Dimension]
The topological dimension $d_t = 4$ remains fixed throughout, with exactly \textbf{1 time dimension} and \textbf{3 to 9 space dimensions}. The spectral dimension $d_s$ increases from 4 to 10 not by adding time dimensions, but by expanding the accessible space dimensions from 3 (external) to 9 (internal). This is consistent with the $Cl(1,9)$ structure: 1 time + 9 space dimensions in the internal space.
\end{remark}

\subsection{Gravitational Potential as Internal Space Energy Gradient}

\begin{hypothesis}[Gravitation as Internal Space Potential]
Gravitation is not a fundamental force but an emergent phenomenon arising from the energy gradient between external and internal spaces:
\begin{equation}
    \Phi(r) = -\frac{GM}{r} = -\tau_0^2 c^2 \cdot f\left(\frac{r}{r_s}\right)
    \label{eq:gravity_as_potential}
\end{equation}
where $f(r/r_s)$ is a dimensionless shape function.
\end{hypothesis}

The gravitational force arises from the tendency of energy to flow "downhill" from the external space (higher potential) to the internal space (lower potential):

\begin{equation}
    \mathbf{F}_g = -\nabla\Phi = -c^2 \tau_0 \nabla\tau
    \label{eq:gravitational_force}
\end{equation}

\begin{remark}[Physical Analogy]
This is analogous to water flowing downhill due to gravitational potential. In our framework, energy flows "inward" due to the internal space potential gradient, and this flow manifests as the gravitational force we observe.
\end{remark}

\subsection{Unified Picture: Two Questions as One Process}

The two seemingly distinct phenomena---dimensional change across the horizon and gravitational potential---are unified in our framework:

\begin{equation}
    \text{Gravitational potential gradient} \to \text{Energy flow inward} \to \text{Dimension expands } 4D \to 10D
    \label{eq:unified_process}
\end{equation}

The torsion field $\tau_{\mu\nu\alpha}$ serves as the bridge:
\begin{enumerate}
    \item It generates gravitational effects (equation \eqref{eq:gravitational_force})
    \item It drives dimensional transition (equation \eqref{eq:spectral_dimension})
\end{enumerate}

\subsection{Rotating Sphere Experiment as Black Hole Analogue}

The rotating sphere experiment discussed in section \ref{sec:experimental} provides a tabletop analogue of black hole physics:

\begin{table}[htbp]
\centering
\caption{Rotating Sphere vs Black Hole Analogy}
\begin{tabular}{@{}ccc@{}}
\toprule
Feature & Rotating Sphere & Black Hole \\
\midrule
Inward mechanism & Centripetal force $m\omega^2 r$ & Gravitational force $GMm/r^2$ \\
Energy flow & External $\to$ Internal space & External $\to$ Internal space \\
Spectral dimension & $d_s^{outer} \to 2$ & $d_s^{outer} \to 2$ \\
Internal dimension & $d_s^{inner} \to 10$ & $d_s^{inner} \to 10$ \\
Nature & Artificial, energy-input required & Natural, self-sustaining \\
\bottomrule
\end{tabular}
\label{tab:sphere_bh_analogy}
\end{table}

The key insight is that both systems exhibit \textbf{inward fractal expansion}---the apparent "contraction" in external coordinates corresponds to expansion in the internal space.

\subsection{Force Direction Differentiation at GUT Scale}

At the Grand Unification scale ($E \sim 10^{16}$ GeV), the three gauge forces (electromagnetic, weak, strong) exhibit a unique behavior:

\begin{theorem}[Unified but Distinguished]
At $E = E_{GUT}$:
\begin{enumerate}
    \item The coupling constants converge: $\alpha_1 \approx \alpha_2 \approx \alpha_3$
    \item The gauge symmetries unify into a single algebraic structure
    \item However, the geometric "directions" (projection angles) remain distinct
\end{enumerate}
\end{theorem}

This leads to a state of \textbf{unified but distinguished} forces:

\begin{equation}
    U(1) \times SU(2) \times SU(3) \xrightarrow{E_{GUT}} G_{unified} \text{ with angular separation}
\end{equation}

The angular separation persists because it reflects the fundamental geometric structure of the internal space:
\begin{itemize}
    \item Electromagnetic: Shallow angle ($15°$) --- closest to external geometry
    \item Weak: Intermediate angle ($30°$)
    \item Strong: Deep angle ($45°$) --- deepest into internal space
\end{itemize}

\begin{remark}[Contrast with Standard GUT]
In standard Grand Unified Theories, the forces become completely indistinguishable at $E_{GUT}$. In our framework, they unify algebraically but remain geometrically distinct through their different "directions" in the internal space.
\end{remark}

\subsection{Summary}

This section establishes:
\begin{enumerate}
    \item The fundamental distinction between topological dimension (fixed) and spectral dimension (variable)
    \item The mechanism of energy state transition across the black hole horizon
    \item The concept of inward fractal expansion resolving the singularity problem
    \item The interpretation of gravity as internal space potential gradient
    \item The unified picture connecting dimensional change and gravitational effects
    \item The rotating sphere experiment as a black hole analogue
    \item The unique behavior of forces at the GUT scale: unified but geometrically distinguished
\end{enumerate}

These insights provide a coherent geometric framework for understanding black hole physics and the unification of fundamental forces.
